%! AP Statistics — Unit 6 Cover Sheet
\documentclass[10pt]{article}
\usepackage{apstats-article}
\usepackage{apstats-boxes}

\begin{document}

% ── Full-bleed navy banner ────────────────────────────────────────────────────
\begin{tikzpicture}[remember picture, overlay]
  \fill[navy] (current page.north west)
    rectangle ([yshift=-1.4in]current page.north east);
\end{tikzpicture}
\vspace{-1.3in}

\begin{center}
  {\color{white}\bfseries\fontsize{28}{34}\selectfont AP Statistics}\\[8pt]
  {\color{goldacc}\bfseries\Large Shepherd \quad 2026--2027}\\[14pt]
  {\color{white}\bfseries\LARGE Unit 6: Inference for Categorical Data: Proportions}\\[4pt]
\end{center}

\vspace{0.30in}

% ── Unit overview ─────────────────────────────────────────────────────────────
\begin{tcolorbox}[colback=sky,colframe=navy,boxrule=0.9pt,arc=2mm,
    left=4mm,right=4mm,top=3mm,bottom=3mm]
  \textbf{Unit Overview.}\quad
  Unit 6 introduces formal statistical inference --- the process of using sample data to draw
  conclusions about unknown population parameters. Students begin with confidence intervals for
  a single population proportion and learn to verify conditions, construct intervals, interpret
  them correctly, and justify claims. They then extend these ideas to significance testing:
  stating hypotheses, calculating test statistics and $p$-values, making decisions, and stating
  conclusions in context. The unit closes by examining potential errors in inference (Type I,
  Type II, power) and by extending all procedures to the two-sample case: confidence intervals
  and significance tests for a difference between two population proportions. This unit establishes
  the reasoning framework applied in every subsequent inference unit.
\end{tcolorbox}

\vspace{0.22in}

% ── Lessons in this unit ──────────────────────────────────────────────────────
\renewcommand{\arraystretch}{1.6}
\begin{skillbox}[Lessons in This Unit]{goldbox}
\begin{tabularx}{\linewidth}{@{} c >{\bfseries}p{0.46\linewidth} X @{}}
  \textbf{\#} & \textbf{Title} & \textbf{Focus} \\
  \hline
  6.1  & Introducing Statistics: Why Be Normal?
       & Distribution shape variation; bridge from sampling distributions to inference. \\
  \hline
  6.2  & Constructing a Confidence Interval for a Population Proportion
       & One-sample $z$-interval; conditions; $SE_{\hat p}$; margin of error; sample size. \\
  \hline
  6.3  & Justifying a Claim Based on a Confidence Interval for a Population Proportion
       & Interpreting CIs; width vs.\ $n$ and $C$; justifying claims. \\
  \hline
  6.4  & Setting Up a Test for a Population Proportion
       & $H_0$ and $H_a$; one-sample $z$-test; checking conditions ($np_0 \ge 10$, etc.). \\
  \hline
  6.5  & Interpreting $p$-Values
       & Test statistic $z$; $p$-value definition and interpretation; null distribution. \\
  \hline
  6.6  & Concluding a Test for a Population Proportion
       & Significance level $\alpha$; reject / fail to reject; conclusion in context. \\
  \hline
  6.7  & Potential Errors When Performing Tests
       & Type I / Type II errors; power; factors affecting error probabilities. \\
  \hline
  6.8  & Confidence Intervals for the Difference of Two Proportions
       & Two-sample $z$-interval for $p_1 - p_2$; conditions; formula and calculation. \\
  \hline
  6.9  & Justifying a Claim Based on a CI for a Difference of Population Proportions
       & Interpreting two-sample CIs; justifying claims about differences. \\
  \hline
  6.10 & Setting Up a Test for the Difference of Two Population Proportions
       & $H_0{:}\;p_1 = p_2$; two-sample $z$-test; conditions with pooled proportion $\hat p_c$. \\
  \hline
  6.11 & Carrying Out a Test for the Difference of Two Population Proportions
       & Pooled $z$-statistic; $p$-value; decision and conclusion in context. \\
  \hline
\end{tabularx}
\end{skillbox}

\vspace{0.18in}

% ── Standards covered ─────────────────────────────────────────────────────────
\begin{skillbox}[AP Statistics Big Ideas \& Learning Objectives --- Unit 6]{greenbox}
\begin{tabularx}{\linewidth}{@{} >{\bfseries}p{0.14\linewidth} X @{}}
  VAR-1.H  & Identify questions suggested by variation in shapes of distributions of samples. \\
  UNC-4.A--E & Identify, verify conditions, calculate, and interpret a one-sample $z$-interval for $p$. \\
  UNC-4.F--H & Interpret CI; justify claim; relate width to $n$ and $C$. \\
  VAR-6.D--F & State $H_0$/$H_a$; identify one-sample $z$-test; verify conditions for testing $p$. \\
  VAR-6.G  & Calculate $z$-statistic and $p$-value for a one-sample proportion test. \\
  DAT-3.A--B & Interpret $p$-value; decide to reject or fail to reject $H_0$; state conclusion. \\
  UNC-5.A--D & Define, calculate, and interpret Type I/II errors and power. \\
  UNC-4.I--L & Identify, verify conditions, and calculate a two-sample $z$-interval for $p_1-p_2$. \\
  UNC-4.M--N & Interpret and justify claims from a two-sample CI for a difference. \\
  VAR-6.H--J & State hypotheses; identify two-sample $z$-test; verify conditions using $\hat p_c$. \\
  VAR-6.K; DAT-3.C--D & Calculate pooled $z$-statistic; interpret $p$-value; conclude in context. \\
\end{tabularx}
\end{skillbox}

\end{document}
